\section*{Hoofdstuk \ref{ch:g12:subduction-magmatism} --- Magmatisme bij subductiezones}

\begin{solution}{exo:g12:subduction-magmatism:1}
(a) een \omterm{def:g12:subduction-magmatism:arc}{vulkanische boog}; (b) \omterm{def:g12:subduction-magmatism:hydrous-melting}{watergeïnduceerd smelten} (fluxsmelten); (c) een \omterm{def:g12:subduction-magmatism:granodiorite}{granodioriet};
(d) een \omterm{def:g12:subduction-magmatism:batholith}{batholiet}; (e) \omterm{def:g12:subduction-magmatism:continental-growth}{continentale groei}.
\end{solution}

\begin{solution}{exo:g12:subduction-magmatism:2}
(1) Minstens $24/0.25 = \qty{96}{km/h}$, dat is $\num{96000}/3600 \approx
\qty{27}{m/s}$. (2) De explosie was minstens $480/96 = 5$ keer sneller.
\end{solution}

\begin{solution}{exo:g12:subduction-magmatism:3}
(1) $838/1214 \approx 0.690$, ongeveer \qty{69}{\percent}. (2) Aan de kant van
de plaat die erover schuift, die niet zinkt. (3) Tussen de trog
en het front ligt de \omterm{def:g11:subduction-zones:slab}{slab} nog niet diep genoeg: hij heeft de diepte nog niet bereikt,
ongeveer honderd kilometer, waar zijn water vrijkomt in een mantel die heet
genoeg is om te smelten zodra hij nat is.
\end{solution}

\begin{solution}{exo:g12:subduction-magmatism:4}
(1) $45/(45 + 15) = 0.75$: \qty{75}{\percent} van de veldspaat is
plagioklaas. (2) Ja: het is grofkorrelig, bevat meer dan
\qty{20}{\percent} kwarts, en zijn aandeel plagioklaas ligt tussen 65 en
\qty{90}{\percent}.
\end{solution}

\begin{solution}{exo:g12:subduction-magmatism:5}
Boogmagma's zijn rijker aan silica, dus stroperiger; en ze dragen
het water dat ze heeft gemaakt, dus zijn ze rijk aan gas. Een hoge viscositeit en een
hoog gasgehalte begunstigen allebei explosieve uitbarstingen.
\end{solution}

\begin{solution}{exo:g12:subduction-magmatism:6}
(1) $H = 150 \times \tan\ang{35} \approx 150 \times 0.700 \approx
\qty{105}{km}$. (2) Het is het wereldgemiddelde.
\end{solution}

\begin{solution}{exo:g12:subduction-magmatism:7}
$d = H/\tan\theta$: $105/\tan\ang{20} \approx 105/0.364 \approx
\qty{288}{km}$ bij \ang{20}, en $105/\tan\ang{45} = \qty{105}{km}$ bij
\ang{45}. Hoe flauwer de \omterm{def:g11:subduction-zones:slab}{slab}, hoe verder van de trog de boog staat.
\end{solution}

\begin{solution}{exo:g12:subduction-magmatism:8}
(1) $T = -5.104 \times 3.16^2 + 132.9 \times 3.16 + 1120.7 \approx -51.0 +
420.0 + 1120.7 \approx \qty{1490}{\celsius}$. (2) Met ongeveer $1490 - 1000 =
\qty{490}{\celsius}$. (3) Droog niet: \qty{1300}{\celsius} ligt onder
\qty{1490}{\celsius}. Nat wel: het ligt boven \qty{1000}{\celsius}.
\end{solution}

\begin{solution}{exo:g12:subduction-magmatism:9}
(1) $2.5/2.7 \approx 0.93$ jaar, ongeveer 340 dagen. (2) De lawine verplaatste
oud gesteente dat al bestond; bogen maken nieuw gesteente uit de mantel. Alleen de
volumes worden vergeleken.
\end{solution}

\begin{solution}{exo:g12:subduction-magmatism:10}
(1) $2.7 \times \num{1000000} = \qty{2.7e6}{km^3}$. (2) $\num{9e9}/2.7
\approx \num{3.3e9}$ jaar, 3.3 miljard jaar. (3) Het is minder dan de ouderdom van
de Aarde, dus bogen hadden de hele \omterm{def:g12:continental-crust:continental-crust}{continentale korst} kunnen maken; maar
\omterm{def:g11:subduction-zones:subduction}{subductie} neemt ook korst terug, dus in feite is veel van de korst
meer dan eens gemaakt en vernietigd.
\end{solution}

\begin{solution}{exo:g12:subduction-magmatism:11}
(1) Ondiepe aardbevingen komen ook voor in de bovenliggende plaat en op de breuk
tussen de twee platen; alleen de diepe liggen in de \omterm{def:g11:subduction-zones:slab}{slab} alleen.
(2) $116 - 80 = \qty{36}{km}$ over \qty{86.4}{km}: $\tan\theta = 36/86.4
\approx 0.417$, $\theta \approx \ang{22.6}$.
\end{solution}

\begin{solution}{exo:g12:subduction-magmatism:12}
In de rechthoekige driehoek van de propositie is de \omterm{def:g11:subduction-zones:slab}{slab} de schuine zijde $L$
en $d$ de zijde die aan de hoek $\theta$ grenst, dus $\cos\theta = d/L$ en $L =
d/\cos\theta$. Voor Japan: $L = 270/\cos\ang{20} \approx 270/0.940 \approx
\qty{287}{km}$. Bij \qty{83}{mm} per jaar: $\num{2.87e8}/83 \approx \num{3.5e6}$
jaar, ongeveer 3.5 miljoen jaar.
\end{solution}

\begin{solution}{exo:g12:subduction-magmatism:13}
(1) $V = \pi r^2 h = \pi \times 0.5^2 \times 8 \approx \qty{6.3}{km^3}$.
(2) Ongeveer $6.3/2.5 \approx 2.5$ keer het volume van de lawine.
(3) $P = \rho g z$: op \qty{6}{km} $2830 \times 9.81 \times 6000 \approx
\qty{1.67e8}{Pa} = \qty{0.17}{GPa}$; op \qty{14}{km} ongeveer
\qty{3.89e8}{Pa} = \qty{0.39}{GPa}.
\end{solution}

\begin{solution}{exo:g12:subduction-magmatism:14}
(1) $2.7 - 3.2 = \qty{-0.5}{km^3}$ per jaar, een verlies van \qty{5e5}{km^3} per
miljoen jaar. (2) $500 \times \num{5e5} = \qty{2.5e8}{km^3}$, dat is
$\num{2.5e8}/\num{9e9} \approx \qty{2.8}{\percent}$ van de korst. (3) Elk
tempo is een schatting, met een onzekerheid die veel groter is dan het verschil van
\qty{0.5}{km^3} per jaar ertussen: binnen de foutmarges houden de twee elkaar in evenwicht.
\end{solution}

\begin{solution}{exo:g12:subduction-magmatism:15}
(1) $\num{20000} \times 10 = \qty{2e5}{km^3}$. (2) $\num{2e5}/1 =
\num{200000}$ jaar voor alle continentale bogen van de wereld samen. Eén
boog krijgt maar een klein deel van dat magma, dus doet één \omterm{def:g12:subduction-magmatism:batholith}{batholiet} er
veel langer over: hij wordt \omterm{def:g11:igneous-rocks:intrusion}{pluton} na \omterm{def:g11:igneous-rocks:intrusion}{pluton} opgebouwd, over miljoenen jaren. (3) Minstens
de \qty{5}{km} gesteente boven zijn top.
\end{solution}

\begin{solution}{pb:g12:subduction-magmatism:1}
\textbf{1.} $111.2 \times \cos\ang{39} \approx 111.2 \times 0.777 \approx
\qty{86.4}{km}$.

\textbf{2.} $(144.1 - 141.0) \times 86.4 \approx 3.1 \times 86.4 \approx
\qty{268}{km}$, ongeveer \qty{270}{km}.

\textbf{3.} $(141.0 - 139.7) \times 86.4 \approx \qty{112}{km}$.

\textbf{4.} Op de bovenliggende plaat, die niet zinkt, ten westen van de
trog: de Pacifische plaat zinkt eronder.

\textbf{5.} Daar ligt de \omterm{def:g11:subduction-zones:slab}{slab} nog niet diep genoeg om zijn water vrij te geven in
een mantel die heet genoeg is om te smelten.

\textbf{6.} $116 - 80 = \qty{36}{km}$.

\textbf{7.} Eén lengtegraad, \qty{86.4}{km}.

\textbf{8.} $\tan\theta = 36/86.4 \approx 0.417$, dus $\theta \approx
\ang{22.6}$.

\textbf{9.} De lijn stijgt \qty{36}{km} per \qty{86.4}{km} naar het oosten, dus
bereikt ze het oppervlak $98/0.417 \approx \qty{235}{km}$ ten oosten van Iwate, op
$141.0 + 235/86.4 \approx \ang{143.7}$~OL: ongeveer $(144.1 - 143.7) \times 86.4
\approx \qty{35}{km}$ ten westen van de trog.

\textbf{10.} De \omterm{def:g11:subduction-zones:slab}{slab} is geen vlak: hij gaat bijna horizontaal de trog in
en buigt steeds sterker omlaag naarmate hij zinkt.

\textbf{11.} $268 \times \tan\ang{20} \approx 268 \times 0.364 \approx
\qty{98}{km}$.

\textbf{12.} Ze komen overeen: beide keren \qty{98}{km}.

\textbf{13.} Beide liggen binnen het bereik voor de wereld, iets ondieper dan het
gemiddelde van \qty{105}{km}.

\textbf{14.} Megata ligt \qty{112}{km} ten westen van Iwate, dus ligt de \omterm{def:g11:subduction-zones:slab}{slab} eronder
ongeveer $98 + 0.417 \times 112 \approx \qty{145}{km}$ diep.

\textbf{15.} In de mantel, ongeveer $98 - 35 = \qty{63}{km}$ onder de basis van
de korst.

\textbf{16.} Het kwam erin dicht bij de rug waar de korst ontstond, uit
zeewater; veel ervan zit in serpentijn.

\textbf{17.} $T = -5.104 \times 3.16^2 + 132.9 \times 3.16 + 1120.7 \approx
\qty{1490}{\celsius}$.

\textbf{18.} Met $1490 - 1000 = 490$ tot $1490 - 950 = \qty{540}{\celsius}$:
ongeveer \qty{500}{\celsius}.

\textbf{19.} De \omterm{def:g11:subduction-zones:slab}{slab} is enkele honderden graden kouder dan de mantel
eromheen, te koud om zelfs met water te smelten: het is de hete mantelwig erboven,
eenmaal nat gemaakt, die smelt.

\textbf{20.} De \omterm{def:g11:subduction-zones:slab}{slab} ligt ongeveer \textbf{\qty{100}{km} onder het vulkanische
front} van het noordoosten van Japan (\qty{98}{km} met beide methoden).
\end{solution}
